\subsection{模的逆向极限}

设 \((\mathbf{A}_{\alpha}, \phi_{\alpha\beta})\) 是环的一个逆向系统（I，§ 10，第 1 号）， \((\mathbf{E}_{\alpha}, f_{\alpha\beta})\) 是交换群（加法记法）的一个逆向系统（I，§10，第1号），并假设每个 \(E_{\alpha}\) 具有一个左 \(A_{\alpha}\)-模结构；此外，假设对 \(\alpha \leqslant \beta\)，\((f_{\alpha\beta}, \phi_{\alpha\beta})\) 是 \(E_{\beta}\) 到 \(E_{\alpha}\) 的一个双态射（§1，第13号），换句话说，

\[
f _ {\alpha \beta} (\lambda_ {\beta} x _ {\beta}) = \phi_ {\alpha \beta} (\lambda_ {\beta}) f _ {\alpha \beta} (x _ {\beta}),\tag{1}
\]

对 \(x_{\beta} \in \mathrm{E}_{\beta}, \lambda_{\beta} \in \mathrm{A}_{\beta}\)；于是由I，§10，第2号可知 \(\mathbf{E} = \varprojlim \mathbf{E}_{\alpha}\) 具有一个关于 \(\mathbf{A} = \varprojlim \mathbf{A}_{\alpha}\) 的左模结构。对所有 \(\alpha \in \mathbf{I}\)，设 \(f: \mathbf{E} \to \mathbf{E}_{\alpha}, \phi_{\alpha}: \mathbf{A} \to \mathbf{A}_{\alpha}\) 为典范映射；则 \((f_{\alpha}, \phi_{\alpha})\) 是 \(\mathbf{E}\) 到 \(\mathbf{E}_{\alpha}\) 的一个双态射。我们将称 \((\mathbf{E}_{\alpha}, f_{\alpha\beta})\) 是左 \(A_{\alpha}\)-模的逆向系统，且A-模E是其逆向极限。

设 \((\mathbf{E}_{\alpha}^{\prime}, f_{\alpha\beta}^{\prime})\) 是左 \(A_{\alpha}\)-模的另一个逆向系统，并且对所有 \(\alpha\)，设 \(u_{\alpha}: E_{\alpha}^{\prime} \to E_{\alpha}\) 为 \(A_{\alpha}\)-线性映射，且这些映射构成一个逆向系统；则 \(u = \lim_{\leftarrow} u_{\alpha}\) 是从 \(\lim_{\leftarrow} E_{\alpha}^{\prime}\) 到 \(\lim_{\leftarrow} E_{\alpha}\) 的一个A-线性映射。

此外：

命题1. 设 \((\mathbf{E}_{\alpha}, f_{\alpha\beta})\) , \((\mathbf{E}_{\alpha}^{\prime}, f_{\alpha\beta}^{\prime})\) , \((\mathbf{E}_{\alpha}^{\prime \prime}, f_{\alpha\beta}^{\prime \prime})\) 是 \(\mathbf{A}_{\alpha}\)-模的三个逆向系统，且 \((u_{\alpha})\) , \((v_{\alpha})\) 是 \(\mathbf{A}_{\alpha}\)-线性映射的两个逆向系统，使得序列

\[
0 \longrightarrow \mathrm{E} _ {\alpha} ^ {\prime} \stackrel {u _ {\alpha}} {\longrightarrow} \mathrm{E} _ {\alpha} \stackrel {v _ {\alpha}} {\longrightarrow} \mathrm{E} _ {\alpha} ^ {\prime \prime}
\]

对所有 \(\alpha\) 正合。于是，记 \(u = \varprojlim u_{\alpha}\) , \(v = \varprojlim v_{\alpha}\)，则序列

\[
0 \longrightarrow \varprojlim \mathrm{E} _ {\alpha} ^ {\prime} \stackrel {u} {\longrightarrow} \varprojlim \mathrm{E} _ {\alpha} \stackrel {v} {\longrightarrow} \varprojlim \mathrm{E} _ {\alpha} ^ {\prime \prime}
\]

是正合的。

由于 \(u_{\alpha}^{-1}(0) = \{0\}\)（对所有 \(\alpha\)），由集合论，III，§7，第2号，命题2可知 \(u^{-1}(0) = \{0\}\)，因此 \(u\) 是单射；进而，诸 \(u_{\alpha}(\mathrm{E}_{\alpha}')\) 构成 \(\mathrm{E}_{\alpha}\) 的子集的一个逆向系统，因此 \(u(\varprojlim \mathrm{E}_{\alpha}') = \varprojlim u_{\alpha}(\mathrm{E}_{\alpha}')\)。由假设 \(u_{\alpha}(\mathrm{E}_{\alpha}') = v_{\alpha}^{-1}(0)\)，故 \(v^{-1}(0) = \varprojlim u_{\alpha}(\mathrm{E}_{\alpha}') = u(\varprojlim \mathrm{E}_{\alpha}')\)（集合论，III，§7，第2号，命题2），证明完毕。

注记. (1) 只需改变记号，命题1及其证明对任意群均成立。

\begin{enumerate}
\def\labelenumi{(\arabic{enumi})}
\setcounter{enumi}{1}
\tightlist
\item
  注意，如果存在正合序列
\end{enumerate}

\[
0 \longrightarrow \mathrm{E} _ {\alpha} ^ {\prime} \stackrel {u _ {\alpha}} {\longrightarrow} \mathrm{E} _ {\alpha} \stackrel {v _ {\alpha}} {\longrightarrow} \mathrm{E} _ {\alpha} ^ {\prime \prime} \longrightarrow 0
\]

并不能由此得出序列

\[
0 \longrightarrow \varprojlim \mathrm{E} _ {\alpha} ^ {\prime} \stackrel {u} {\longrightarrow} \varprojlim \mathrm{E} _ {\alpha} \stackrel {v} {\longrightarrow} \varprojlim \mathrm{E} _ {\alpha} ^ {\prime \prime} \longrightarrow 0
\]

是正合的；换言之，满射线性映射的逆向系统的逆向极限不一定是满射（参见练习1）。

假设现在诸 \(\mathbf{A}_{\alpha}\) 等于同一个环 \(\mathbf{A}\) 且 \(\phi_{\alpha \beta}\) 等于 \(1_{\mathbf{A}}\)；则对A-模的每个逆向系统 \((\mathrm{E}_{\alpha}, f_{\alpha \beta})\)，\(\mathbf{E} = \varprojlim \mathbf{E}_{\alpha}\) 是一个A-模。设 \(\mathbf{F}\) 是一个A-模，且对所有 \(\alpha\)，设 \(u_{\alpha}: \mathbf{F} \to \mathbf{E}_{\alpha}\) 是A-线性映射，使得 \((u_{\alpha})\) 是映射的一个逆向系统；则 \(u = \varprojlim u_{\alpha}\) 是从 \(\mathbf{F}\) 到 \(\varprojlim \mathbf{E}_{\alpha}\) 的一个A-线性映射。反之，对每个A-线性映射 \(v: \mathbf{F} \to \varprojlim \mathbf{E}_{\alpha}\)，族 \(v_{\alpha} = f_{\alpha} \circ v\) 是A-线性映射的一个逆向系统，使得 \(v = \varprojlim v_{\alpha}\)。另一方面，我们注意到对 \(\alpha \leqslant \beta\) 映射

\[
\operatorname{Hom} \left(\mathrm{l} _ {\mathrm{F}}, f _ {\alpha \beta}\right) = \bar {f} _ {\alpha \beta}: \operatorname{Hom} _ {\mathrm{A}} (\mathrm{F}, \mathrm{E} _ {\beta}) \rightarrow \operatorname{Hom} _ {\mathrm{A}} (\mathrm{F}, \mathrm{E} _ {\alpha})
\]

是一个Z-模同态，使得 \((\mathrm{Hom}_{\mathbf{A}}(\mathrm{F}, \mathrm{E}_{\alpha}), \vec{f}_{\alpha\beta})\) 是Z-模的一个逆向系统；由于 \(\vec{f}_{\alpha\beta}(v_{\beta}) = f_{\alpha\beta} \circ v_{\beta}\)，上述论述因此可表述如下：

命题2. 对由A-模组成的每个逆向系统 \((\mathbf{E}_{\alpha}, f_{\alpha \beta})\) 以及每个A-模F，典范映射 \(u \mapsto (f_{\alpha} \circ u)\) 是Z-模同构

\[
l _ {\mathrm{F}}: \operatorname{Hom} _ {\mathrm{A}} (\mathrm{F}, \varprojlim \mathrm{E} _ {\alpha}) \rightarrow \varprojlim \operatorname{Hom} _ {\mathrm{A}} (\mathrm{F}, \mathrm{E} _ {\alpha}).\tag{2}
\]

推论. 对于每个A-模同态 \(v: \mathbf{F} \to \mathbf{F}'\) ，

\[
\bar {v} _ {\alpha} = \operatorname{Hom} (v, 1 _ {\mathrm{E} _ {\alpha}}): \operatorname{Hom} (\mathrm{F} ^ {\prime}, \mathrm{E} _ {\alpha}) \rightarrow \operatorname{Hom} (\mathrm{F}, \mathrm{E} _ {\alpha})
\]

构成Z-线性映射的一个逆向系统，且图表

\[
\begin{array}{c} \operatorname{Hom} (\mathrm{F} ^ {\prime}, \varprojlim \mathrm{E} _ {\alpha}) \xrightarrow {l _ {\mathrm{F} ^ {\prime}}} \varprojlim \operatorname{Hom} (\mathrm{F} ^ {\prime}, \mathrm{E} _ {\alpha}) \\ \operatorname{Hom} (v, 1 _ {\mathrm{E}}) \Biggl \downarrow \qquad \qquad \qquad \qquad \qquad \qquad \qquad \qquad \qquad \qquad \qquad \qquad \qquad \qquad \qquad \qquad \qquad \qquad \qquad \qquad \qquad \qquad \qquad \qquad \qquad \qquad \qquad \qquad \qquad \qquad \qquad \qquad \qquad \qquad \\ \operatorname{Hom} (\mathrm{F}, \varprojlim \mathrm{E} _ {\alpha}) \xrightarrow {l _ {\mathrm{F}}} \varprojlim \operatorname{Hom} (\mathrm{F}, \mathrm{E} _ {\alpha}) \end{array}\tag{3}
\]

是交换的。

对所有 \(u \in \operatorname{Hom}(\mathbf{F}', \varprojlim \mathbf{E}_{\alpha})\)，由定义有 \(l_{\mathbf{F}}(u \circ v) = (f_{\alpha} \circ u \circ v)\)，且图(3)的交换性由定义直接得出。

